\section{Theoretical Framework: Perez and the Phases of Technological Revolutions}

Carlota Perez's model of techno-economic paradigms divides each technological revolution into two broad phases: \textbf{installation} and \textbf{deployment}~\cite{perez2002}. During installation, financial capital floods into the new technology, building infrastructure and generating speculative bubbles. This phase is characterized by tensions between paper economy and real economy, between potential supply and effective demand, and between rich and poor~\cite{palma2009}. The transition to deployment requires the disciplining of finance capital and the rise of productive capital, often through institutional reform, crisis, or both~\cite{swilling2019}.

The installation period is not merely a prelude. It is the phase in which the infrastructure for the next paradigm is built---often wastefully, often through bubbles, but built nonetheless. The railways of the 1840s, the electrification of the 1890s, and the fiber-optic buildout of the late 1990s all followed this pattern: overinvestment, crash, and then the productive use of the infrastructure that remained~\cite{perez2002}.

Applied to AI, the installation period corresponds roughly to the period from 2020 to the present. Financial capital has driven massive investment in data centers, GPUs, and training runs. The ``bubble'' is not necessarily in AI company valuations alone; it is in the physical infrastructure---the data center capacity, the chip fabrication, the energy contracts. This infrastructure will remain after the speculative fever subsides. The question is who will control it and on what terms.

Perez's framework is used here as an analytical lens, not as proof that AI has entered a new historical phase. It generates questions about the relationship between infrastructure investment, innovation, and institutional change. If AI is moving from an installation-oriented regime toward broader deployment, one would expect changes in four observable domains: investment priorities (from frontier training toward inference and application deployment), innovation patterns (from scaling toward efficiency, reliability, and integration), institutional arrangements (the emergence of evaluation regimes, standards, and certification), and the distribution of returns (who controls infrastructure, who captures productivity gains).

Each of these four indicators can be operationalized, even coarsely. Investment priorities can be tracked through disclosed training-run costs (Epoch AI's database) and through capital-expenditure guidance from hyperscalers. Innovation patterns can be tracked through the composition of arXiv submissions (efficiency and inference-time-computation papers versus scaling papers) and through model-release composition (small efficient models versus frontier-scale models). Institutional arrangements can be tracked through the number of jurisdictions with binding AI evaluation regimes and through the number of labs with published evaluator-access agreements. Distribution of returns can be tracked through the share of inference revenue captured by the top three labs and through the share of compute capacity controlled by the top three providers. None of these indicators is currently measured consistently across the four dimensions; specifying them here is a precondition for the empirical work the framework calls for, not a substitute for it.

These are proposed indicators, not validated predictions. The connection between each indicator and the theoretical framework must be justified, and the framework's value lies in organizing observations into a coherent account of technological change---and in making the interpretation vulnerable to evidence that contradicts it. The strongest counter-evidence is the continued improvement in benchmark performance and the persistence of frontier capability gaps; if the deployment shift were not occurring, the rentier turn would not be predicted. The indicators are not yet measured consistently over time, and the theory currently organizes the narrative more than it independently explains the observed events. This is acknowledged as a limitation.

Perez developed her framework around railways, electrification, automobiles, and IT. AI has different cost structures: near-zero marginal cost of inference, data network effects, and open-weight competition. Two ways AI breaks the Perez model are worth naming. First, the installation phase in AI is not driven by a single general-purpose technology but by a stack (chips, models, cloud, power), each with its own investment cycle. Second, the deployment phase may not produce a ``golden age'' of broadly shared gains, because the infrastructure is proprietary from the outset---unlike the railways, where the track was a common carrier, or the early internet, where the protocol layer could not be captured. These differences do not invalidate the framework; they specify where it must be supplemented by attention to institutional design.
